% === Begin Label Data ===
% ID: 726
% 难度星级: 
% 标签: 
% 备注: 
% 组卷引用次数: 0
% === End  Label Data ===

\begin{problem}{2013}{G}{湖南卷（文）}{15}{数列}
对于 $ E=\{a_1,a_2,\cdots,a_{100}\} $ 的子集 $ X=\{a_{i_1},a_{i_2},\cdots,a_{i_k}\} $, 定义 $ X $ 的“特征数列”为 $ x_1,x_2,\cdots,x_{100} $, 其中 $ x_{i_1}=x_{i_2}=\cdots=x_{i_k}=1 $, 其余项均为 $ 0 $, 例如: 子集 $ \{a_2,a_3\} $ 的“特征数列”为 $ 0 $, $ 1 $, $ 0 $, $ \cdots $, $ 0 $.

（1）子集 $ \{a_1,a_3,a_5\} $ 的“特征数列”的前 $ 3 $ 项和等于 \underline{\hspace{4em}}.

（2）若 $ E $ 的子集 $ P $ 的“特征数列” $ p_1,p_2,\cdots,p_{100} $ 满足 $ p_1=1 $, $ p_i+p_{i+1}=1 $, $ 1\leq i\leq 99 $; $ E $ 的子集 $ Q $ 的“特征数列” $ q_1,q_2,\cdots,q_{100} $ 满足 $ q_1=1 $, $ q_j+q_{j+1}+q_{j+2}=1 $, $ 1\leq j\leq 98 $, 则 $ P\cap Q $ 的元素个数为 \underline{\hspace{4em}}.
\end{problem}

\begin{answer}

\end{answer}

\begin{solutions}

\end{solutions}